\subsection{ARCH-136-DEGREE-TWO: concentrating the vertical defect in one degree}
\label{sec:arch136degreetwo}

\paragraph{Archive locators.}
The remaining vertical condition is on cote 136, physical p.~117,
footer 116, \texttt{ARCHIVE\_136\_scan117.tex}; horizontal uniformity
is on physical p.~119, footer 118, \texttt{ARCHIVE\_136\_scan119.tex}.
The common-source question is physical p.~121, footer 120,
\texttt{ARCHIVE\_136\_scan121.tex}.

\paragraph{Statement.}
For every finite nonempty family
$C_1,\ldots,C_t\in\mathcal A_{\mathbf Z}$ there is a simplicial
strict bigraded differential PD $S=\Gamma_{\mathbf Z}(T)$-algebra
$A$ with strict maps to every $C_j$ and the following properties:
\[
 \begin{gathered}
 A_0=S,\qquad A_n^{0,*}=\mathbf Z[0],\qquad
 H^*((A_n)_{\mathbf Q})=S_{\mathbf Q},\\
 H_r(NA_*^{w,q})=0\quad(w\geq\mu(q),\ r\geq0),\\
 H^0(A_n^{w,*})=S^w,\quad H^1(A_n^{w,*})=0
                                      \quad(w\geq W_1),\\
 H^q(A_n^{w,*})=0\quad(q\geq3,\ w\geq W_q).
 \end{gathered}                                           \tag{136.DT1}
\]
All cutoffs are independent of the simplex dimension and the
construction stage.  The sole missing vertical admissibility
condition is eventual vanishing of $H^2(A_n^{w,*})$, whose groups
are integer torsion.  No finite-prime or bounded-exponent assertion
for these degree-two groups is made: the repairs can move defects
from arbitrarily high degrees down to degree two.

\paragraph{Inputs and distinction.}
The proof uses the accepted AQ-136-FIRST-SKELETON replacement,
AQ-136-RELATIVE-H1, and the horizontal attachment from
AQ-136-PARTIAL-SOURCE.  It supplies a new strict simplicial higher
repair that preserves simplex zero and all the given target maps.
It does not claim arbitrary relative completion: that would
contradict AQ-136-HIGHER-COMPLETION.  Instead it isolates exactly
the degree at which the present repair no longer works.

\paragraph{Proof.}
First apply AQ-136-FIRST-SKELETON to the finite family.  This gives
admissible replacements $C'_j\to C_j$ with $(C'_j)_0=S$ exactly.
It suffices to construct the source for the replacements, then
compose.  Write $C_j$ for them for the rest of the proof.  Choose
once and for all the horizontal cutoffs $\mu(q)\geq1$, a common
cutoff $W_1\geq1$ for their degree-zero coefficient comparison
and degree-one vanishing, and, for each $q\geq3$, a cutoff
$W_q\geq1$ such that
\[
 H^q((C_j)_n^{w,*})=H^{q-1}((C_j)_n^{w,*})=0
                \quad(w\geq W_q,\ \text{all }n,j).         \tag{136.DT2}
\]
These are finite maxima of the targets' admissibility bounds.

We describe a repair $A\to V(A)$ under any current stage satisfying
$A_0=S$, weight zero equal to $\mathbf Z[0]$, and exact rational
vertical cohomology in every weight.  For each $q\geq3$, $w\geq W_q$,
simplex dimension $n$, and data
\[
 z\in A_n^{w,q},\quad dz=0,\quad
 b\in A_n^{w,q-1},\quad db=N z,
                          \qquad N\in\mathbf Z_{>0},     \tag{136.DT3}
\]
make a separate attachment.  There is a set of all these data.
Let
\[
       M_n=\mathbf Z[\Delta[n]]/
                       \mathbf Z[\operatorname{sk}_0\Delta[n]]
\]
be the reduced representable simplicial module.  It is free in
every simplicial degree and has degree-zero term zero.  A simplex
$z$ or $b$ in (136.DT3) defines a simplicial module map from the
representable into the corresponding row of $A$.  Its restriction
to every vertex is zero because both form degrees are positive
and $A_0=S$.  Hence these maps factor through $M_n$.

Adjoin free generators $c(M_n)$ and $t(M_n)$ in form degrees
$q-1$ and $q-2$, both of weight $w$, with differential
\[
                   dc=z,\qquad dt=Nc-b.                 \tag{136.DT4}
\]
Here $z,b$ denote the just-defined linear maps from $M_n$.
Thus the formulas are actual simplicial linear assignments, not
independent choices in unrelated simplex dimensions.  Equation
$db=Nz$ gives $d^2=0$.  Extend by the strict graded derivation and
special divided-power rules to the free graded PD extension.
These rules remain consistent on even divided powers because the
differential of an even element is odd and squares to zero.

Every prescribed map $f_j:A\to C_j$ extends across this attachment.
At simplex $n$, choose $c_j$ of form degree $q-1$ with
$dc_j=f_j(z)$ using (136.DT2).  The cochain $Nc_j-f_j(b)$ is closed,
so choose $t_j$ of degree $q-2$ with
$dt_j=Nc_j-f_j(b)$, again by (136.DT2).
Both degrees are strictly positive because $q\geq3$.
Consequently both chosen cochains restrict to zero on vertices,
since $(C_j)_0=S$.  They define maps out of the same $M_n$ and
satisfy (136.DT4).  The free universal property supplies a strict
simplicial target map.  Different copies of the attachment need
not have compatible choices: they are free over their common old
stage, and any new relations among them will be treated by later
repairs.  This construction addresses that duplication directly
without claiming the one-Moore-pair exactness range for a family.

The attachment leaves $A_0=S$ and weight zero unchanged.  In every
fixed simplex the module of new generators is free, and over
$\mathbf Q$ the change $c'=c-b/N$ gives
\[
                       dc'=0,\qquad dt=Nc'.
\]
It is an invertible change of free graded PD generators.  Thus,
rationally and simplexwise, the extension is a tensor product
with free vertical unit disks.  They contract onto their constants,
so the extension preserves rational vertical cohomology in every
weight.  This argument applies to a simultaneous set of attachments
by passing through finite subsets.  The same observation proves
preservation after their filtered union.

Each old high-weight class in a degree $q\geq3$ is killed by this
step.  Indeed rational vertical exactness says that its cohomology
class is integer torsion, so there are $N>0$ and $b$ with $db=Nz$.
One can see this directly from localization of the cohomology
module, without requiring torsion-free underlying terms.  The
corresponding attachment in (136.DT3) has $dc=z$.

Start from $S$ and repeatedly apply three kinds of step:
\begin{enumerate}
 \item the horizontal free-cell step, at the fixed cutoffs $\mu(q)$;
 \item the higher repair $V$ just constructed, at the fixed $W_q$;
 \item the accepted relative repair of AQ-136-RELATIVE-H1, at $W_1$.
\end{enumerate}
Each step preserves $A_0=S$, weight zero, all target maps, and exact
rational vertical cohomology.  The last assertion for the third
step is stronger: its rationalization is the same actual algebra,
by the cited theorem.  Its hypotheses therefore remain satisfied
at every iteration.  The horizontal step needs only the target
row acyclicity, and the new higher step needs only (136.DT2).
Hence every operation is always available with its original fixed
bounds.

Take the filtered colimit after countably many rounds.  Exactness
of filtered colimits of modules and of normalization proves all
the desired vanishings: every old horizontal class in its indicated
weight tail is killed at the next horizontal step; every old
vertical class in degree $q\geq3$ and weight at least $W_q$ is killed
at the next higher step; every old degree-one class is killed at
the next relative $H^1$ step.  The latter also sends every old
zero-degree class in its range to a coefficient.  Coefficients
remain injective because $A_0=S$ and degeneracies from simplex
zero have face retractions.  The colimit therefore has exactly
the claimed degree-zero coefficient cohomology.  Every individual
class in a filtered colimit comes from a stage cohomology class,
so the possibility that later steps create new classes does not
invalidate this argument: those classes are killed in their own
subsequent round.

The rational transitions are quasi-isomorphisms onto the coefficient
cohomology, so the rational assertion also passes to the colimit.
It follows that its remaining degree-two cohomology is integer
torsion.  No step claims a uniform annihilator or prime support
for that degree.  This proves (136.DT1).\hfill$\square$

\paragraph{Why the same proof stops at degree two.}
If $q=2$, the new $t$ in (136.DT4) has degree zero.  A chosen target
primitive $t_j$ restricts on vertices to possibly different
coefficient values.  Its differential vanishes there, but the
cochain itself need not vanish.  Therefore it need not factor
through the reduced representable $M_n$; arbitrary choices cannot
be made while retaining $A_0=S$ and the same maps to every target.
The obstruction is the vector of these vertex coefficients, modulo
changes induced by the choice of the primitive $c_j$ and by adding
a constant coefficient to $t_j$.  In a two-vertex situation this
is precisely the secondary coefficient residue behind
AQ-136-HIGHER-COMPLETION and AQ-136-DERIVED-COMPATIBILITY.

Thus the unresolved part of this construction is a fixed-form-degree
relative problem.  Resolving those coefficient residues, or replacing
the source/maps so that they agree in a uniform weight tail, would
complete this construction.  The theorem neither establishes that
agreement nor refutes integral connectedness.  It also does not
retain ARCH-136-DISK-SUPPORT's finite-prime bound: the price of
concentrating the defects is that degree two can now involve
arbitrarily many primes.
